% =============================================================
%  MAT231-F27 Calculus II — Lecture 25
%  Unit 25: Systems of linear ODEs via eigenvalues
%
%  Build from the repository root:   npm run materials -- lecture-25
%  This writes lecture-25.pdf (slides, with overlays) and
%  lecture-25-handout.pdf (print version) into this folder.
% =============================================================
\documentclass[10pt]{wosm-slides}

\usepackage{pgfplots}
\pgfplotsset{compat=1.18}

\newcommand{\R}{\mathbb{R}}
\newcommand{\vx}{\mathbf{x}}
% \vv is also defined by newtxmath, which wosm.sty loads with Source Serif;
% \DeclareRobustCommand (not \newcommand) works with or without it.
\DeclareRobustCommand{\vv}{\mathbf{v}}
\newcommand{\vc}{\mathbf{c}}

\course{MAT231-F27}{Calculus II}
\title[L25 \textperiodcentered\ Linear systems]{Systems of linear differential equations}
\subtitle{Eigenvalues, phase portraits, and the matrix exponential}
\author{Dr.\ Ali Khudiyev}
\institute{WoSM Academy}
\date{Lecture 25\enspace\textperiodcentered\enspace Unit 25\enspace\textperiodcentered\enspace Fall 2027}

% A small phase portrait: the box [-2,2]^2 with axes.
\newcommand{\phasebox}{%
  \draw[wosm-rule] (-2,0) -- (2,0) (0,-2) -- (0,2);
  \draw[wosm-muted] (-2,-2) rectangle (2,2);}

\begin{document}

\maketitle

% -------------------------------------------------------------
\begin{frame}{Last time}
  \begin{itemize}
    \item $ay'' + by' + cy = 0$: the characteristic equation $ar^2 + br + c = 0$; real, double and complex
      roots; Euler's formula.
    \item Free and damped oscillations; forced oscillations and resonance.
  \end{itemize}
  \vskip0.6em
  Today: several unknown functions at once, $\vx' = A\vx$. The characteristic equation becomes the
  eigenvalue equation $\det(A - \lambda I) = 0$ --- and Linear Algebra and Calculus meet.
\end{frame}

\begin{frame}{This lecture}
  \tableofcontents
  \vskip1em
  {\small\color{wosm-muted}
  Reading: Apostol, \emph{Calculus, Volume 2}, Chapter 7.\par}
\end{frame}

% =============================================================
\section{Linear systems}

\begin{frame}{Systems of differential equations}
  A \textbf{linear system with constant coefficients}: \ $\vx'(t) = A\, \vx(t)$, with $\vx(t) \in \R^n$ and
  $A$ a constant $n \times n$ matrix.
  \pause
  \begin{example}[Two rooms]
    Two rooms lose heat to the outside, and exchange heat with each other, at rate $1$. With $x_1$, $x_2$
    their temperatures above the outside temperature:
    \[ x_1' = -x_1 + (x_2 - x_1), \quad x_2' = -x_2 + (x_1 - x_2), \qquad
       A = \begin{pmatrix} -2 & 1 \\ 1 & -2 \end{pmatrix} . \]
  \end{example}
\end{frame}

\begin{frame}{A second-order equation is a system}
  For $y'' + by' + cy = 0$, put $x_1 = y$ and $x_2 = y'$. Then $x_1' = x_2$ and $x_2' = -c x_1 - b x_2$:
  \[ \vx' = \begin{pmatrix} 0 & 1 \\ -c & -b \end{pmatrix} \vx , \qquad
     \det(A - \lambda I) = \lambda^2 + b\lambda + c , \]
  the characteristic polynomial of Lecture 24.
  \pause
  \begin{itemize}
    \item In the same way, any equation of order $n$ becomes a first-order system in $\R^n$.
    \item Numerical solvers (Euler's method and its refinements) work with first-order systems, so
      this is how higher-order equations are solved on a computer.
  \end{itemize}
\end{frame}

\begin{frame}{Existence, uniqueness, superposition}
  \begin{theorem}
    For every $\vx_0 \in \R^n$, $\vx' = A\vx$ has exactly one solution with $\vx(0) = \vx_0$, and it is
    defined for all $t$.
  \end{theorem}
  \pause
  \begin{itemize}
    \item Sums and multiples of solutions are solutions, and $\vx \mapsto \vx(0)$ is a linear bijection
      onto $\R^n$: the solutions form an $n$-dimensional vector space.
    \item So we need $n$ independent solutions; every solution is a combination of them.
    \item Geometrically, solutions are the flow lines of the vector field $\vx \mapsto A\vx$
      (Lecture 22). Their pictures in the $(x_1, x_2)$-plane form the \textbf{phase portrait}.
  \end{itemize}
\end{frame}

% =============================================================
\section{Solutions from eigenvectors}

\begin{frame}{Straight-line solutions}
  Try $\vx(t) = e^{\lambda t}\, \vv$ with a constant vector $\vv \ne \mathbf 0$:
  \[ \vx' = \lambda e^{\lambda t} \vv, \qquad A\vx = e^{\lambda t} A\vv . \]
  \pause
  So $e^{\lambda t}\vv$ is a solution exactly when $A\vv = \lambda\vv$: \ $\vv$ is an \textbf{eigenvector}
  and $\lambda$ an \textbf{eigenvalue} of $A$ (Linear Algebra, Unit 17).
  \pause
  \begin{itemize}
    \item The solution moves along the line through $\vv$: outwards if $\lambda > 0$, inwards if $\lambda < 0$.
    \item The eigenvalues are the roots of $\det(A - \lambda I) = 0$, the analogue of the characteristic
      equation.
  \end{itemize}
\end{frame}

\begin{frame}{The general solution}
  \begin{theorem}
    If $A$ has $n$ linearly independent eigenvectors $\vv_1, \dots, \vv_n$ with eigenvalues
    $\lambda_1, \dots, \lambda_n$, every solution of $\vx' = A\vx$ is
    \[ \vx(t) = c_1 e^{\lambda_1 t} \vv_1 + \dots + c_n e^{\lambda_n t} \vv_n , \]
    and the $c_i$ are found from $\vx(0) = c_1 \vv_1 + \dots + c_n \vv_n$.
  \end{theorem}
  \pause
  \textbf{Why}: with $V = (\vv_1 \cdots \vv_n)$, $A = V\Lambda V^{-1}$ (Linear Algebra, Unit 18). Put
  $\vx = V\mathbf y$: then $\mathbf y' = \Lambda\mathbf y$, that is $y_i' = \lambda_i y_i$ for each $i$ ---
  $n$ separate equations of Lecture 23, with $y_i = c_i e^{\lambda_i t}$.
  \pause
  \vskip0.4em
  {\small Eigenvectors \textbf{decouple} the system. Each eigendirection evolves on its own.\par}
\end{frame}

\begin{frame}{Example: two rooms}
  $A = \begin{pmatrix} -2 & 1 \\ 1 & -2 \end{pmatrix}$: \ $\det(A - \lambda I) = (\lambda + 2)^2 - 1$, so
  $\lambda = -1$ with $\vv = (1, 1)$, and $\lambda = -3$ with $\vv = (1, -1)$.
  \pause
  \vskip0.4em
  Start with room 1 warm, room 2 not: $\vx(0) = (1, 0) = \frac12 (1, 1) + \frac12 (1, -1)$. So
  \[ \vx(t) = \frac12 e^{-t} \begin{pmatrix} 1 \\ 1 \end{pmatrix} + \frac12 e^{-3t} \begin{pmatrix} 1 \\ -1 \end{pmatrix} . \]
  \pause
  \begin{itemize}
    \item The \textbf{difference} $x_1 - x_2 = e^{-3t}$ disappears fast: the rooms equalise.
    \item The \textbf{average} $\frac12(x_1 + x_2) = \frac12 e^{-t}$ decays slowly: both cool together.
  \end{itemize}
\end{frame}

\begin{frame}{Complex eigenvalues}
  If $\lambda = \alpha + i\beta$ ($\beta \ne 0$) is an eigenvalue of a real matrix $A$, with eigenvector $\vv$,
  the real and imaginary parts of $e^{\lambda t}\vv$ are two real solutions (as in Lecture 24).
  \pause
  \vskip0.4em
  \textbf{Example.} $A = \begin{pmatrix} -1 & 2 \\ -2 & -1 \end{pmatrix}$: \ $\det(A - \lambda I) = (\lambda + 1)^2 + 4$,
  so $\lambda = -1 \pm 2i$. \ For $\lambda = -1 + 2i$, \ $\vv = (1, i)$:
  \[ e^{\lambda t}\vv = e^{-t}(\cos 2t + i \sin 2t) \begin{pmatrix} 1 \\ i \end{pmatrix}
     = e^{-t}\begin{pmatrix} \cos 2t \\ -\sin 2t \end{pmatrix} + i\, e^{-t}\begin{pmatrix} \sin 2t \\ \cos 2t \end{pmatrix} . \]
  \pause
  The solution with $\vx(0) = (1, 0)$ is $e^{-t}(\cos 2t, -\sin 2t)$: it turns clockwise and shrinks ---
  a \textbf{spiral}.
\end{frame}

% =============================================================
\section{Phase portraits and stability}

\begin{frame}{Phase portraits}
  \begin{columns}[t]
    \begin{column}{0.32\textwidth}
      \centering
      \begin{tikzpicture}[scale=0.72]
        \phasebox
        \draw[wosm-accent] plot coordinates {(1.440,-1.360) (1.401,-1.313) (1.363,-1.267) (1.328,-1.222) (1.294,-1.177) (1.261,-1.133) (1.231,-1.090) (1.202,-1.047) (1.175,-1.005) (1.145,-0.955) (1.117,-0.905) (1.093,-0.855) (1.071,-0.805) (1.052,-0.755) (1.037,-0.704) (1.025,-0.652) (1.016,-0.598) (1.011,-0.544) (1.010,-0.487) (1.014,-0.429) (1.020,-0.378) (1.030,-0.324) (1.044,-0.269) (1.062,-0.210) (1.080,-0.161) (1.100,-0.110) (1.124,-0.057) (1.151,-0.001) (1.182,0.059) (1.207,0.105) (1.235,0.153) (1.265,0.203) (1.298,0.256) (1.333,0.311) (1.371,0.368) (1.412,0.427) (1.441,0.469) (1.472,0.511) (1.504,0.555) (1.537,0.600) (1.572,0.647) (1.609,0.695) (1.647,0.745) (1.687,0.796) (1.729,0.850) (1.773,0.904) (1.819,0.961) (1.867,1.020) (1.917,1.080) (1.969,1.143)};
        \draw[wosm-accent, -latex] (1.100,-0.110) -- (1.124,-0.057);
        \draw[wosm-accent] plot coordinates {(1.600,-1.200) (1.581,-1.150) (1.564,-1.099) (1.549,-1.048) (1.537,-0.997) (1.526,-0.944) (1.519,-0.891) (1.513,-0.837) (1.511,-0.781) (1.511,-0.724) (1.514,-0.666) (1.520,-0.606) (1.529,-0.544) (1.542,-0.480) (1.554,-0.431) (1.568,-0.380) (1.584,-0.327) (1.603,-0.273) (1.624,-0.217) (1.648,-0.159) (1.674,-0.099) (1.703,-0.037) (1.735,0.027) (1.770,0.094) (1.796,0.140) (1.822,0.188) (1.851,0.236) (1.881,0.286) (1.912,0.338) (1.945,0.391) (1.980,0.445)};
        \draw[wosm-accent, -latex] (1.568,-0.380) -- (1.584,-0.327);
        \draw[wosm-accent] plot coordinates {(1.900,-0.900) (1.903,-0.845) (1.908,-0.789) (1.915,-0.731) (1.925,-0.672) (1.937,-0.612) (1.952,-0.550) (1.969,-0.486) (1.989,-0.420)};
        \draw[wosm-accent, -latex] (1.925,-0.672) -- (1.937,-0.612);
        \draw[wosm-accent] plot coordinates {(1.360,-1.440) (1.313,-1.401) (1.267,-1.363) (1.222,-1.328) (1.177,-1.294) (1.133,-1.261) (1.090,-1.231) (1.047,-1.202) (1.005,-1.175) (0.955,-1.145) (0.905,-1.117) (0.855,-1.093) (0.805,-1.071) (0.755,-1.052) (0.704,-1.037) (0.652,-1.025) (0.598,-1.016) (0.544,-1.011) (0.487,-1.010) (0.429,-1.014) (0.378,-1.020) (0.324,-1.030) (0.269,-1.044) (0.210,-1.062) (0.161,-1.080) (0.110,-1.100) (0.057,-1.124) (0.001,-1.151) (-0.059,-1.182) (-0.105,-1.207) (-0.153,-1.235) (-0.203,-1.265) (-0.256,-1.298) (-0.311,-1.333) (-0.368,-1.371) (-0.427,-1.412) (-0.469,-1.441) (-0.511,-1.472) (-0.555,-1.504) (-0.600,-1.537) (-0.647,-1.572) (-0.695,-1.609) (-0.745,-1.647) (-0.796,-1.687) (-0.850,-1.729) (-0.904,-1.773) (-0.961,-1.819) (-1.020,-1.867) (-1.080,-1.917) (-1.143,-1.969)};
        \draw[wosm-accent, -latex] (0.110,-1.100) -- (0.057,-1.124);
        \draw[wosm-accent] plot coordinates {(1.200,-1.600) (1.150,-1.581) (1.099,-1.564) (1.048,-1.549) (0.997,-1.537) (0.944,-1.526) (0.891,-1.519) (0.837,-1.513) (0.781,-1.511) (0.724,-1.511) (0.666,-1.514) (0.606,-1.520) (0.544,-1.529) (0.480,-1.542) (0.431,-1.554) (0.380,-1.568) (0.327,-1.584) (0.273,-1.603) (0.217,-1.624) (0.159,-1.648) (0.099,-1.674) (0.037,-1.703) (-0.027,-1.735) (-0.094,-1.770) (-0.140,-1.796) (-0.188,-1.822) (-0.236,-1.851) (-0.286,-1.881) (-0.338,-1.912) (-0.391,-1.945) (-0.445,-1.980)};
        \draw[wosm-accent, -latex] (0.380,-1.568) -- (0.327,-1.584);
        \draw[wosm-accent] plot coordinates {(0.900,-1.900) (0.845,-1.903) (0.789,-1.908) (0.731,-1.915) (0.672,-1.925) (0.612,-1.937) (0.550,-1.952) (0.486,-1.969) (0.420,-1.989)};
        \draw[wosm-accent, -latex] (0.672,-1.925) -- (0.612,-1.937);
        \draw[wosm-accent] plot coordinates {(-1.360,1.440) (-1.313,1.401) (-1.267,1.363) (-1.222,1.328) (-1.177,1.294) (-1.133,1.261) (-1.090,1.231) (-1.047,1.202) (-1.005,1.175) (-0.955,1.145) (-0.905,1.117) (-0.855,1.093) (-0.805,1.071) (-0.755,1.052) (-0.704,1.037) (-0.652,1.025) (-0.598,1.016) (-0.544,1.011) (-0.487,1.010) (-0.429,1.014) (-0.378,1.020) (-0.324,1.030) (-0.269,1.044) (-0.210,1.062) (-0.161,1.080) (-0.110,1.100) (-0.057,1.124) (-0.001,1.151) (0.059,1.182) (0.105,1.207) (0.153,1.235) (0.203,1.265) (0.256,1.298) (0.311,1.333) (0.368,1.371) (0.427,1.412) (0.469,1.441) (0.511,1.472) (0.555,1.504) (0.600,1.537) (0.647,1.572) (0.695,1.609) (0.745,1.647) (0.796,1.687) (0.850,1.729) (0.904,1.773) (0.961,1.819) (1.020,1.867) (1.080,1.917) (1.143,1.969)};
        \draw[wosm-accent, -latex] (-0.110,1.100) -- (-0.057,1.124);
        \draw[wosm-accent] plot coordinates {(-1.200,1.600) (-1.150,1.581) (-1.099,1.564) (-1.048,1.549) (-0.997,1.537) (-0.944,1.526) (-0.891,1.519) (-0.837,1.513) (-0.781,1.511) (-0.724,1.511) (-0.666,1.514) (-0.606,1.520) (-0.544,1.529) (-0.480,1.542) (-0.431,1.554) (-0.380,1.568) (-0.327,1.584) (-0.273,1.603) (-0.217,1.624) (-0.159,1.648) (-0.099,1.674) (-0.037,1.703) (0.027,1.735) (0.094,1.770) (0.140,1.796) (0.188,1.822) (0.236,1.851) (0.286,1.881) (0.338,1.912) (0.391,1.945) (0.445,1.980)};
        \draw[wosm-accent, -latex] (-0.380,1.568) -- (-0.327,1.584);
        \draw[wosm-accent] plot coordinates {(-0.900,1.900) (-0.845,1.903) (-0.789,1.908) (-0.731,1.915) (-0.672,1.925) (-0.612,1.937) (-0.550,1.952) (-0.486,1.969) (-0.420,1.989)};
        \draw[wosm-accent, -latex] (-0.672,1.925) -- (-0.612,1.937);
        \draw[wosm-accent] plot coordinates {(-1.440,1.360) (-1.401,1.313) (-1.363,1.267) (-1.328,1.222) (-1.294,1.177) (-1.261,1.133) (-1.231,1.090) (-1.202,1.047) (-1.175,1.005) (-1.145,0.955) (-1.117,0.905) (-1.093,0.855) (-1.071,0.805) (-1.052,0.755) (-1.037,0.704) (-1.025,0.652) (-1.016,0.598) (-1.011,0.544) (-1.010,0.487) (-1.014,0.429) (-1.020,0.378) (-1.030,0.324) (-1.044,0.269) (-1.062,0.210) (-1.080,0.161) (-1.100,0.110) (-1.124,0.057) (-1.151,0.001) (-1.182,-0.059) (-1.207,-0.105) (-1.235,-0.153) (-1.265,-0.203) (-1.298,-0.256) (-1.333,-0.311) (-1.371,-0.368) (-1.412,-0.427) (-1.441,-0.469) (-1.472,-0.511) (-1.504,-0.555) (-1.537,-0.600) (-1.572,-0.647) (-1.609,-0.695) (-1.647,-0.745) (-1.687,-0.796) (-1.729,-0.850) (-1.773,-0.904) (-1.819,-0.961) (-1.867,-1.020) (-1.917,-1.080) (-1.969,-1.143)};
        \draw[wosm-accent, -latex] (-1.100,0.110) -- (-1.124,0.057);
        \draw[wosm-accent] plot coordinates {(-1.600,1.200) (-1.581,1.150) (-1.564,1.099) (-1.549,1.048) (-1.537,0.997) (-1.526,0.944) (-1.519,0.891) (-1.513,0.837) (-1.511,0.781) (-1.511,0.724) (-1.514,0.666) (-1.520,0.606) (-1.529,0.544) (-1.542,0.480) (-1.554,0.431) (-1.568,0.380) (-1.584,0.327) (-1.603,0.273) (-1.624,0.217) (-1.648,0.159) (-1.674,0.099) (-1.703,0.037) (-1.735,-0.027) (-1.770,-0.094) (-1.796,-0.140) (-1.822,-0.188) (-1.851,-0.236) (-1.881,-0.286) (-1.912,-0.338) (-1.945,-0.391) (-1.980,-0.445)};
        \draw[wosm-accent, -latex] (-1.568,0.380) -- (-1.584,0.327);
        \draw[wosm-accent] plot coordinates {(-1.900,0.900) (-1.903,0.845) (-1.908,0.789) (-1.915,0.731) (-1.925,0.672) (-1.937,0.612) (-1.952,0.550) (-1.969,0.486) (-1.989,0.420)};
        \draw[wosm-accent, -latex] (-1.925,0.672) -- (-1.937,0.612);
        \draw[wosm-brick, densely dashed] (-2,-2) -- (2,2) (-2,2) -- (2,-2);
      \end{tikzpicture}
      \par\small \textbf{Saddle}: \\ $\lambda = 3, -1$. \\ $A = \begin{pmatrix} 1 & 2 \\ 2 & 1 \end{pmatrix}$
    \end{column}
    \begin{column}{0.32\textwidth}
      \centering
      \begin{tikzpicture}[scale=0.72]
        \phasebox
        \draw[wosm-accent] plot coordinates {(-2.000,2.000) (-1.912,1.912) (-1.827,1.827) (-1.747,1.747) (-1.670,1.670) (-1.596,1.596) (-1.526,1.526) (-1.458,1.458) (-1.394,1.394) (-1.333,1.333) (-1.274,1.274) (-1.218,1.218) (-1.164,1.164) (-1.113,1.113) (-1.064,1.064) (-1.017,1.017) (-0.972,0.972) (-0.929,0.929) (-0.888,0.888) (-0.849,0.849) (-0.811,0.811) (-0.776,0.776) (-0.709,0.709) (-0.647,0.647) (-0.592,0.592) (-0.541,0.541) (-0.494,0.494) (-0.451,0.451) (-0.412,0.412) (-0.377,0.377) (-0.329,0.329) (-0.287,0.287) (-0.251,0.251) (-0.210,0.210) (-0.167,0.167) (-0.128,0.128) (-0.089,0.089) (-0.052,0.052) (-0.016,0.016)};
        \draw[wosm-accent, -latex] (-0.849,0.849) -- (-0.811,0.811);
        \draw[wosm-accent] plot coordinates {(2.000,2.000) (1.970,1.970) (1.912,1.912) (1.855,1.855) (1.800,1.800) (1.747,1.747) (1.695,1.695) (1.645,1.645) (1.596,1.596) (1.549,1.549) (1.503,1.503) (1.458,1.458) (1.415,1.415) (1.373,1.373) (1.333,1.333) (1.293,1.293) (1.255,1.255) (1.218,1.218) (1.182,1.182) (1.129,1.129) (1.080,1.080) (1.032,1.032) (0.986,0.986) (0.943,0.943) (0.901,0.901) (0.862,0.862) (0.824,0.824) (0.787,0.787) (0.741,0.741) (0.698,0.698) (0.657,0.657) (0.619,0.619) (0.583,0.583) (0.541,0.541) (0.501,0.501) (0.465,0.465) (0.425,0.425) (0.388,0.388) (0.350,0.350) (0.310,0.310) (0.271,0.271) (0.233,0.233) (0.197,0.197) (0.160,0.160) (0.124,0.124) (0.088,0.088) (0.052,0.052) (0.016,0.016)};
        \draw[wosm-accent, -latex] (0.901,0.901) -- (0.862,0.862);
        \draw[wosm-accent] plot coordinates {(2.000,-2.000) (1.912,-1.912) (1.827,-1.827) (1.747,-1.747) (1.670,-1.670) (1.596,-1.596) (1.526,-1.526) (1.458,-1.458) (1.394,-1.394) (1.333,-1.333) (1.274,-1.274) (1.218,-1.218) (1.164,-1.164) (1.113,-1.113) (1.064,-1.064) (1.017,-1.017) (0.972,-0.972) (0.929,-0.929) (0.888,-0.888) (0.849,-0.849) (0.811,-0.811) (0.776,-0.776) (0.709,-0.709) (0.647,-0.647) (0.592,-0.592) (0.541,-0.541) (0.494,-0.494) (0.451,-0.451) (0.412,-0.412) (0.377,-0.377) (0.329,-0.329) (0.287,-0.287) (0.251,-0.251) (0.210,-0.210) (0.167,-0.167) (0.128,-0.128) (0.089,-0.089) (0.052,-0.052) (0.016,-0.016)};
        \draw[wosm-accent, -latex] (0.849,-0.849) -- (0.811,-0.811);
        \draw[wosm-accent] plot coordinates {(-2.000,-2.000) (-1.970,-1.970) (-1.912,-1.912) (-1.855,-1.855) (-1.800,-1.800) (-1.747,-1.747) (-1.695,-1.695) (-1.645,-1.645) (-1.596,-1.596) (-1.549,-1.549) (-1.503,-1.503) (-1.458,-1.458) (-1.415,-1.415) (-1.373,-1.373) (-1.333,-1.333) (-1.293,-1.293) (-1.255,-1.255) (-1.218,-1.218) (-1.182,-1.182) (-1.129,-1.129) (-1.080,-1.080) (-1.032,-1.032) (-0.986,-0.986) (-0.943,-0.943) (-0.901,-0.901) (-0.862,-0.862) (-0.824,-0.824) (-0.787,-0.787) (-0.741,-0.741) (-0.698,-0.698) (-0.657,-0.657) (-0.619,-0.619) (-0.583,-0.583) (-0.541,-0.541) (-0.501,-0.501) (-0.465,-0.465) (-0.425,-0.425) (-0.388,-0.388) (-0.350,-0.350) (-0.310,-0.310) (-0.271,-0.271) (-0.233,-0.233) (-0.197,-0.197) (-0.160,-0.160) (-0.124,-0.124) (-0.088,-0.088) (-0.052,-0.052) (-0.016,-0.016)};
        \draw[wosm-accent, -latex] (-0.901,-0.901) -- (-0.862,-0.862);
        \draw[wosm-accent] plot coordinates {(-1.000,2.000) (-0.941,1.926) (-0.885,1.856) (-0.832,1.788) (-0.782,1.723) (-0.733,1.661) (-0.687,1.601) (-0.644,1.544) (-0.602,1.489) (-0.563,1.436) (-0.525,1.386) (-0.489,1.337) (-0.455,1.290) (-0.423,1.246) (-0.393,1.203) (-0.363,1.161) (-0.309,1.084) (-0.261,1.012) (-0.217,0.946) (-0.178,0.885) (-0.142,0.829) (-0.111,0.777) (-0.082,0.729) (-0.057,0.684) (-0.024,0.623) (0.004,0.569) (0.027,0.521) (0.052,0.464) (0.071,0.415) (0.088,0.363) (0.101,0.311) (0.110,0.256) (0.111,0.204) (0.103,0.152) (0.084,0.104) (0.057,0.062) (0.023,0.024)};
        \draw[wosm-accent, -latex] (-0.217,0.946) -- (-0.178,0.885);
        \draw[wosm-accent] plot coordinates {(2.000,1.000) (1.956,1.000) (1.871,0.997) (1.790,0.992) (1.715,0.986) (1.643,0.977) (1.576,0.967) (1.512,0.956) (1.451,0.943) (1.394,0.929) (1.339,0.915) (1.288,0.900) (1.239,0.884) (1.169,0.860) (1.105,0.835) (1.045,0.809) (0.989,0.783) (0.937,0.757) (0.888,0.731) (0.843,0.705) (0.786,0.671) (0.734,0.638) (0.686,0.606) (0.642,0.575) (0.601,0.545) (0.554,0.509) (0.511,0.475) (0.471,0.443) (0.429,0.407) (0.390,0.374) (0.350,0.338) (0.309,0.300) (0.269,0.263) (0.231,0.227) (0.192,0.190) (0.155,0.154) (0.118,0.118) (0.082,0.082) (0.047,0.046)};
        \draw[wosm-accent, -latex] (0.843,0.705) -- (0.786,0.671);
        \draw[wosm-accent] plot coordinates {(1.000,-2.000) (0.941,-1.926) (0.885,-1.856) (0.832,-1.788) (0.782,-1.723) (0.733,-1.661) (0.687,-1.601) (0.644,-1.544) (0.602,-1.489) (0.563,-1.436) (0.525,-1.386) (0.489,-1.337) (0.455,-1.290) (0.423,-1.246) (0.393,-1.203) (0.363,-1.161) (0.309,-1.084) (0.261,-1.012) (0.217,-0.946) (0.178,-0.885) (0.142,-0.829) (0.111,-0.777) (0.082,-0.729) (0.057,-0.684) (0.024,-0.623) (-0.004,-0.569) (-0.027,-0.521) (-0.052,-0.464) (-0.071,-0.415) (-0.088,-0.363) (-0.101,-0.311) (-0.110,-0.256) (-0.111,-0.204) (-0.103,-0.152) (-0.084,-0.104) (-0.057,-0.062) (-0.023,-0.024)};
        \draw[wosm-accent, -latex] (0.217,-0.946) -- (0.178,-0.885);
        \draw[wosm-accent] plot coordinates {(-2.000,-1.000) (-1.956,-1.000) (-1.871,-0.997) (-1.790,-0.992) (-1.715,-0.986) (-1.643,-0.977) (-1.576,-0.967) (-1.512,-0.956) (-1.451,-0.943) (-1.394,-0.929) (-1.339,-0.915) (-1.288,-0.900) (-1.239,-0.884) (-1.169,-0.860) (-1.105,-0.835) (-1.045,-0.809) (-0.989,-0.783) (-0.937,-0.757) (-0.888,-0.731) (-0.843,-0.705) (-0.786,-0.671) (-0.734,-0.638) (-0.686,-0.606) (-0.642,-0.575) (-0.601,-0.545) (-0.554,-0.509) (-0.511,-0.475) (-0.471,-0.443) (-0.429,-0.407) (-0.390,-0.374) (-0.350,-0.338) (-0.309,-0.300) (-0.269,-0.263) (-0.231,-0.227) (-0.192,-0.190) (-0.155,-0.154) (-0.118,-0.118) (-0.082,-0.082) (-0.047,-0.046)};
        \draw[wosm-accent, -latex] (-0.843,-0.705) -- (-0.786,-0.671);
        \draw[wosm-accent] plot coordinates {(0.000,2.000) (0.029,1.941) (0.057,1.884) (0.082,1.829) (0.107,1.777) (0.130,1.726) (0.151,1.677) (0.171,1.629) (0.207,1.540) (0.239,1.456) (0.266,1.379) (0.290,1.306) (0.310,1.239) (0.327,1.176) (0.341,1.117) (0.353,1.062) (0.363,1.010) (0.374,0.939) (0.380,0.874) (0.384,0.815) (0.385,0.762) (0.383,0.697) (0.377,0.640) (0.369,0.588) (0.356,0.531) (0.342,0.482) (0.323,0.430) (0.300,0.378) (0.277,0.333) (0.251,0.290) (0.220,0.245) (0.189,0.205) (0.156,0.164) (0.122,0.126) (0.087,0.088) (0.052,0.052) (0.016,0.016)};
        \draw[wosm-accent, -latex] (0.380,0.874) -- (0.384,0.815);
        \draw[wosm-accent] plot coordinates {(2.000,0.000) (1.941,0.029) (1.884,0.057) (1.829,0.082) (1.777,0.107) (1.726,0.130) (1.677,0.151) (1.629,0.171) (1.540,0.207) (1.456,0.239) (1.379,0.266) (1.306,0.290) (1.239,0.310) (1.176,0.327) (1.117,0.341) (1.062,0.353) (1.010,0.363) (0.939,0.374) (0.874,0.380) (0.815,0.384) (0.762,0.385) (0.697,0.383) (0.640,0.377) (0.588,0.369) (0.531,0.356) (0.482,0.342) (0.430,0.323) (0.378,0.300) (0.333,0.277) (0.290,0.251) (0.245,0.220) (0.205,0.189) (0.164,0.156) (0.126,0.122) (0.088,0.087) (0.052,0.052) (0.016,0.016)};
        \draw[wosm-accent, -latex] (0.874,0.380) -- (0.815,0.384);
        \draw[wosm-accent] plot coordinates {(0.000,-2.000) (-0.029,-1.941) (-0.057,-1.884) (-0.082,-1.829) (-0.107,-1.777) (-0.130,-1.726) (-0.151,-1.677) (-0.171,-1.629) (-0.207,-1.540) (-0.239,-1.456) (-0.266,-1.379) (-0.290,-1.306) (-0.310,-1.239) (-0.327,-1.176) (-0.341,-1.117) (-0.353,-1.062) (-0.363,-1.010) (-0.374,-0.939) (-0.380,-0.874) (-0.384,-0.815) (-0.385,-0.762) (-0.383,-0.697) (-0.377,-0.640) (-0.369,-0.588) (-0.356,-0.531) (-0.342,-0.482) (-0.323,-0.430) (-0.300,-0.378) (-0.277,-0.333) (-0.251,-0.290) (-0.220,-0.245) (-0.189,-0.205) (-0.156,-0.164) (-0.122,-0.126) (-0.087,-0.088) (-0.052,-0.052) (-0.016,-0.016)};
        \draw[wosm-accent, -latex] (-0.380,-0.874) -- (-0.384,-0.815);
        \draw[wosm-accent] plot coordinates {(-2.000,0.000) (-1.941,-0.029) (-1.884,-0.057) (-1.829,-0.082) (-1.777,-0.107) (-1.726,-0.130) (-1.677,-0.151) (-1.629,-0.171) (-1.540,-0.207) (-1.456,-0.239) (-1.379,-0.266) (-1.306,-0.290) (-1.239,-0.310) (-1.176,-0.327) (-1.117,-0.341) (-1.062,-0.353) (-1.010,-0.363) (-0.939,-0.374) (-0.874,-0.380) (-0.815,-0.384) (-0.762,-0.385) (-0.697,-0.383) (-0.640,-0.377) (-0.588,-0.369) (-0.531,-0.356) (-0.482,-0.342) (-0.430,-0.323) (-0.378,-0.300) (-0.333,-0.277) (-0.290,-0.251) (-0.245,-0.220) (-0.205,-0.189) (-0.164,-0.156) (-0.126,-0.122) (-0.088,-0.087) (-0.052,-0.052) (-0.016,-0.016)};
        \draw[wosm-accent, -latex] (-0.874,-0.380) -- (-0.815,-0.384);
        \draw[wosm-accent] plot coordinates {(1.000,2.000) (1.000,1.956) (0.997,1.871) (0.992,1.790) (0.986,1.715) (0.977,1.643) (0.967,1.576) (0.956,1.512) (0.943,1.451) (0.929,1.394) (0.915,1.339) (0.900,1.288) (0.884,1.239) (0.860,1.169) (0.835,1.105) (0.809,1.045) (0.783,0.989) (0.757,0.937) (0.731,0.888) (0.705,0.843) (0.671,0.786) (0.638,0.734) (0.606,0.686) (0.575,0.642) (0.545,0.601) (0.509,0.554) (0.475,0.511) (0.443,0.471) (0.407,0.429) (0.374,0.390) (0.338,0.350) (0.300,0.309) (0.263,0.269) (0.227,0.231) (0.190,0.192) (0.154,0.155) (0.118,0.118) (0.082,0.082) (0.046,0.047)};
        \draw[wosm-accent, -latex] (0.705,0.843) -- (0.671,0.786);
        \draw[wosm-accent] plot coordinates {(2.000,-1.000) (1.926,-0.941) (1.856,-0.885) (1.788,-0.832) (1.723,-0.782) (1.661,-0.733) (1.601,-0.687) (1.544,-0.644) (1.489,-0.602) (1.436,-0.563) (1.386,-0.525) (1.337,-0.489) (1.290,-0.455) (1.246,-0.423) (1.203,-0.393) (1.161,-0.363) (1.084,-0.309) (1.012,-0.261) (0.946,-0.217) (0.885,-0.178) (0.829,-0.142) (0.777,-0.111) (0.729,-0.082) (0.684,-0.057) (0.623,-0.024) (0.569,0.004) (0.521,0.027) (0.464,0.052) (0.415,0.071) (0.363,0.088) (0.311,0.101) (0.256,0.110) (0.204,0.111) (0.152,0.103) (0.104,0.084) (0.062,0.057) (0.024,0.023)};
        \draw[wosm-accent, -latex] (0.946,-0.217) -- (0.885,-0.178);
        \draw[wosm-accent] plot coordinates {(-1.000,-2.000) (-1.000,-1.956) (-0.997,-1.871) (-0.992,-1.790) (-0.986,-1.715) (-0.977,-1.643) (-0.967,-1.576) (-0.956,-1.512) (-0.943,-1.451) (-0.929,-1.394) (-0.915,-1.339) (-0.900,-1.288) (-0.884,-1.239) (-0.860,-1.169) (-0.835,-1.105) (-0.809,-1.045) (-0.783,-0.989) (-0.757,-0.937) (-0.731,-0.888) (-0.705,-0.843) (-0.671,-0.786) (-0.638,-0.734) (-0.606,-0.686) (-0.575,-0.642) (-0.545,-0.601) (-0.509,-0.554) (-0.475,-0.511) (-0.443,-0.471) (-0.407,-0.429) (-0.374,-0.390) (-0.338,-0.350) (-0.300,-0.309) (-0.263,-0.269) (-0.227,-0.231) (-0.190,-0.192) (-0.154,-0.155) (-0.118,-0.118) (-0.082,-0.082) (-0.046,-0.047)};
        \draw[wosm-accent, -latex] (-0.705,-0.843) -- (-0.671,-0.786);
        \draw[wosm-accent] plot coordinates {(-2.000,1.000) (-1.926,0.941) (-1.856,0.885) (-1.788,0.832) (-1.723,0.782) (-1.661,0.733) (-1.601,0.687) (-1.544,0.644) (-1.489,0.602) (-1.436,0.563) (-1.386,0.525) (-1.337,0.489) (-1.290,0.455) (-1.246,0.423) (-1.203,0.393) (-1.161,0.363) (-1.084,0.309) (-1.012,0.261) (-0.946,0.217) (-0.885,0.178) (-0.829,0.142) (-0.777,0.111) (-0.729,0.082) (-0.684,0.057) (-0.623,0.024) (-0.569,-0.004) (-0.521,-0.027) (-0.464,-0.052) (-0.415,-0.071) (-0.363,-0.088) (-0.311,-0.101) (-0.256,-0.110) (-0.204,-0.111) (-0.152,-0.103) (-0.104,-0.084) (-0.062,-0.057) (-0.024,-0.023)};
        \draw[wosm-accent, -latex] (-0.946,0.217) -- (-0.885,0.178);
        \draw[wosm-brick, densely dashed] (-2,-2) -- (2,2) (-2,2) -- (2,-2);
      \end{tikzpicture}
      \par\small \textbf{Stable node}: \\ $\lambda = -1, -3$. \\ $A = \begin{pmatrix} -2 & 1 \\ 1 & -2 \end{pmatrix}$
    \end{column}
    \begin{column}{0.32\textwidth}
      \centering
      \begin{tikzpicture}[scale=0.72]
        \phasebox
        \begin{scope}
          \clip (-2,-2) rectangle (2,2);
        \draw[wosm-accent] plot coordinates {(2.000,0.000) (1.991,-0.033) (1.973,-0.099) (1.952,-0.163) (1.929,-0.226) (1.905,-0.288) (1.878,-0.349) (1.850,-0.408) (1.820,-0.465) (1.789,-0.521) (1.756,-0.576) (1.721,-0.629) (1.685,-0.680) (1.648,-0.730) (1.609,-0.778) (1.569,-0.824) (1.528,-0.869) (1.485,-0.912) (1.442,-0.953) (1.398,-0.992) (1.353,-1.029) (1.307,-1.065) (1.260,-1.099) (1.212,-1.131) (1.164,-1.161) (1.115,-1.189) (1.066,-1.215) (1.017,-1.240) (0.967,-1.263) (0.916,-1.283) (0.866,-1.302) (0.815,-1.319) (0.764,-1.335) (0.713,-1.348) (0.662,-1.360) (0.611,-1.370) (0.561,-1.378) (0.510,-1.384) (0.460,-1.389) (0.410,-1.392) (0.335,-1.393) (0.262,-1.390) (0.190,-1.384) (0.119,-1.375) (0.049,-1.362) (-0.019,-1.346) (-0.085,-1.326) (-0.149,-1.304) (-0.212,-1.279) (-0.272,-1.251) (-0.330,-1.220) (-0.386,-1.187) (-0.439,-1.152) (-0.490,-1.115) (-0.538,-1.075) (-0.584,-1.034) (-0.627,-0.991) (-0.668,-0.946) (-0.705,-0.900) (-0.740,-0.853) (-0.772,-0.805) (-0.801,-0.756) (-0.828,-0.706) (-0.851,-0.656) (-0.872,-0.605) (-0.890,-0.553) (-0.905,-0.502) (-0.918,-0.450) (-0.927,-0.399) (-0.934,-0.347) (-0.939,-0.296) (-0.940,-0.246) (-0.939,-0.180) (-0.933,-0.115) (-0.923,-0.051) (-0.909,0.010) (-0.892,0.069) (-0.870,0.127) (-0.846,0.181) (-0.818,0.233) (-0.787,0.283) (-0.754,0.329) (-0.718,0.372) (-0.681,0.412) (-0.641,0.449) (-0.599,0.483) (-0.556,0.513) (-0.512,0.540) (-0.467,0.564) (-0.410,0.588) (-0.352,0.607) (-0.294,0.622) (-0.236,0.631) (-0.179,0.635) (-0.123,0.634) (-0.069,0.629) (-0.016,0.619) (0.035,0.606) (0.084,0.588) (0.130,0.567) (0.181,0.538) (0.228,0.504) (0.271,0.467) (0.308,0.427) (0.341,0.384) (0.368,0.340) (0.390,0.294) (0.410,0.239) (0.422,0.184) (0.428,0.130) (0.428,0.077) (0.421,0.026) (0.409,-0.023) (0.390,-0.074) (0.364,-0.122) (0.333,-0.163) (0.298,-0.199) (0.255,-0.233) (0.210,-0.259) (0.162,-0.276) (0.110,-0.287) (0.058,-0.289) (0.004,-0.282) (-0.045,-0.265) (-0.092,-0.239) (-0.131,-0.205) (-0.163,-0.163) (-0.184,-0.117) (-0.195,-0.068) (-0.193,-0.016) (-0.178,0.033) (-0.151,0.076) (-0.113,0.109) (-0.065,0.128) (-0.014,0.131) (0.035,0.114) (0.072,0.078) (0.089,0.030) (0.079,-0.020) (0.042,-0.054) (-0.008,-0.056) (-0.040,-0.017) (-0.014,0.026)};
        \draw[wosm-accent, -latex] (-0.940,-0.246) -- (-0.939,-0.180);
        \draw[wosm-accent] plot coordinates {(-2.000,-0.000) (-1.991,0.033) (-1.973,0.099) (-1.952,0.163) (-1.929,0.226) (-1.905,0.288) (-1.878,0.349) (-1.850,0.408) (-1.820,0.465) (-1.789,0.521) (-1.756,0.576) (-1.721,0.629) (-1.685,0.680) (-1.648,0.730) (-1.609,0.778) (-1.569,0.824) (-1.528,0.869) (-1.485,0.912) (-1.442,0.953) (-1.398,0.992) (-1.353,1.029) (-1.307,1.065) (-1.260,1.099) (-1.212,1.131) (-1.164,1.161) (-1.115,1.189) (-1.066,1.215) (-1.017,1.240) (-0.967,1.263) (-0.916,1.283) (-0.866,1.302) (-0.815,1.319) (-0.764,1.335) (-0.713,1.348) (-0.662,1.360) (-0.611,1.370) (-0.561,1.378) (-0.510,1.384) (-0.460,1.389) (-0.410,1.392) (-0.335,1.393) (-0.262,1.390) (-0.190,1.384) (-0.119,1.375) (-0.049,1.362) (0.019,1.346) (0.085,1.326) (0.149,1.304) (0.212,1.279) (0.272,1.251) (0.330,1.220) (0.386,1.187) (0.439,1.152) (0.490,1.115) (0.538,1.075) (0.584,1.034) (0.627,0.991) (0.668,0.946) (0.705,0.900) (0.740,0.853) (0.772,0.805) (0.801,0.756) (0.828,0.706) (0.851,0.656) (0.872,0.605) (0.890,0.553) (0.905,0.502) (0.918,0.450) (0.927,0.399) (0.934,0.347) (0.939,0.296) (0.940,0.246) (0.939,0.180) (0.933,0.115) (0.923,0.051) (0.909,-0.010) (0.892,-0.069) (0.870,-0.127) (0.846,-0.181) (0.818,-0.233) (0.787,-0.283) (0.754,-0.329) (0.718,-0.372) (0.681,-0.412) (0.641,-0.449) (0.599,-0.483) (0.556,-0.513) (0.512,-0.540) (0.467,-0.564) (0.410,-0.588) (0.352,-0.607) (0.294,-0.622) (0.236,-0.631) (0.179,-0.635) (0.123,-0.634) (0.069,-0.629) (0.016,-0.619) (-0.035,-0.606) (-0.084,-0.588) (-0.130,-0.567) (-0.181,-0.538) (-0.228,-0.504) (-0.271,-0.467) (-0.308,-0.427) (-0.341,-0.384) (-0.368,-0.340) (-0.390,-0.294) (-0.410,-0.239) (-0.422,-0.184) (-0.428,-0.130) (-0.428,-0.077) (-0.421,-0.026) (-0.409,0.023) (-0.390,0.074) (-0.364,0.122) (-0.333,0.163) (-0.298,0.199) (-0.255,0.233) (-0.210,0.259) (-0.162,0.276) (-0.110,0.287) (-0.058,0.289) (-0.004,0.282) (0.045,0.265) (0.092,0.239) (0.131,0.205) (0.163,0.163) (0.184,0.117) (0.195,0.068) (0.193,0.016) (0.178,-0.033) (0.151,-0.076) (0.113,-0.109) (0.065,-0.128) (0.014,-0.131) (-0.035,-0.114) (-0.072,-0.078) (-0.089,-0.030) (-0.079,0.020) (-0.042,0.054) (0.008,0.056) (0.040,0.017) (0.014,-0.026)};
        \draw[wosm-accent, -latex] (0.940,0.246) -- (0.939,0.180);
        \draw[wosm-accent] plot coordinates {(0.000,2.000) (0.033,1.991) (0.099,1.973) (0.163,1.952) (0.226,1.929) (0.288,1.905) (0.349,1.878) (0.408,1.850) (0.465,1.820) (0.521,1.789) (0.576,1.756) (0.629,1.721) (0.680,1.685) (0.730,1.648) (0.778,1.609) (0.824,1.569) (0.869,1.528) (0.912,1.485) (0.953,1.442) (0.992,1.398) (1.029,1.353) (1.065,1.307) (1.099,1.260) (1.131,1.212) (1.161,1.164) (1.189,1.115) (1.215,1.066) (1.240,1.017) (1.263,0.967) (1.283,0.916) (1.302,0.866) (1.319,0.815) (1.335,0.764) (1.348,0.713) (1.360,0.662) (1.370,0.611) (1.378,0.561) (1.384,0.510) (1.389,0.460) (1.392,0.410) (1.393,0.335) (1.390,0.262) (1.384,0.190) (1.375,0.119) (1.362,0.049) (1.346,-0.019) (1.326,-0.085) (1.304,-0.149) (1.279,-0.212) (1.251,-0.272) (1.220,-0.330) (1.187,-0.386) (1.152,-0.439) (1.115,-0.490) (1.075,-0.538) (1.034,-0.584) (0.991,-0.627) (0.946,-0.668) (0.900,-0.705) (0.853,-0.740) (0.805,-0.772) (0.756,-0.801) (0.706,-0.828) (0.656,-0.851) (0.605,-0.872) (0.553,-0.890) (0.502,-0.905) (0.450,-0.918) (0.399,-0.927) (0.347,-0.934) (0.296,-0.939) (0.246,-0.940) (0.180,-0.939) (0.115,-0.933) (0.051,-0.923) (-0.010,-0.909) (-0.069,-0.892) (-0.127,-0.870) (-0.181,-0.846) (-0.233,-0.818) (-0.283,-0.787) (-0.329,-0.754) (-0.372,-0.718) (-0.412,-0.681) (-0.449,-0.641) (-0.483,-0.599) (-0.513,-0.556) (-0.540,-0.512) (-0.564,-0.467) (-0.588,-0.410) (-0.607,-0.352) (-0.622,-0.294) (-0.631,-0.236) (-0.635,-0.179) (-0.634,-0.123) (-0.629,-0.069) (-0.619,-0.016) (-0.606,0.035) (-0.588,0.084) (-0.567,0.130) (-0.538,0.181) (-0.504,0.228) (-0.467,0.271) (-0.427,0.308) (-0.384,0.341) (-0.340,0.368) (-0.294,0.390) (-0.239,0.410) (-0.184,0.422) (-0.130,0.428) (-0.077,0.428) (-0.026,0.421) (0.023,0.409) (0.074,0.390) (0.122,0.364) (0.163,0.333) (0.199,0.298) (0.233,0.255) (0.259,0.210) (0.276,0.162) (0.287,0.110) (0.289,0.058) (0.282,0.004) (0.265,-0.045) (0.239,-0.092) (0.205,-0.131) (0.163,-0.163) (0.117,-0.184) (0.068,-0.195) (0.016,-0.193) (-0.033,-0.178) (-0.076,-0.151) (-0.109,-0.113) (-0.128,-0.065) (-0.131,-0.014) (-0.114,0.035) (-0.078,0.072) (-0.030,0.089) (0.020,0.079) (0.054,0.042) (0.056,-0.008) (0.017,-0.040) (-0.026,-0.014)};
        \draw[wosm-accent, -latex] (0.246,-0.940) -- (0.180,-0.939);
        \draw[wosm-accent] plot coordinates {(0.000,-2.000) (-0.033,-1.991) (-0.099,-1.973) (-0.163,-1.952) (-0.226,-1.929) (-0.288,-1.905) (-0.349,-1.878) (-0.408,-1.850) (-0.465,-1.820) (-0.521,-1.789) (-0.576,-1.756) (-0.629,-1.721) (-0.680,-1.685) (-0.730,-1.648) (-0.778,-1.609) (-0.824,-1.569) (-0.869,-1.528) (-0.912,-1.485) (-0.953,-1.442) (-0.992,-1.398) (-1.029,-1.353) (-1.065,-1.307) (-1.099,-1.260) (-1.131,-1.212) (-1.161,-1.164) (-1.189,-1.115) (-1.215,-1.066) (-1.240,-1.017) (-1.263,-0.967) (-1.283,-0.916) (-1.302,-0.866) (-1.319,-0.815) (-1.335,-0.764) (-1.348,-0.713) (-1.360,-0.662) (-1.370,-0.611) (-1.378,-0.561) (-1.384,-0.510) (-1.389,-0.460) (-1.392,-0.410) (-1.393,-0.335) (-1.390,-0.262) (-1.384,-0.190) (-1.375,-0.119) (-1.362,-0.049) (-1.346,0.019) (-1.326,0.085) (-1.304,0.149) (-1.279,0.212) (-1.251,0.272) (-1.220,0.330) (-1.187,0.386) (-1.152,0.439) (-1.115,0.490) (-1.075,0.538) (-1.034,0.584) (-0.991,0.627) (-0.946,0.668) (-0.900,0.705) (-0.853,0.740) (-0.805,0.772) (-0.756,0.801) (-0.706,0.828) (-0.656,0.851) (-0.605,0.872) (-0.553,0.890) (-0.502,0.905) (-0.450,0.918) (-0.399,0.927) (-0.347,0.934) (-0.296,0.939) (-0.246,0.940) (-0.180,0.939) (-0.115,0.933) (-0.051,0.923) (0.010,0.909) (0.069,0.892) (0.127,0.870) (0.181,0.846) (0.233,0.818) (0.283,0.787) (0.329,0.754) (0.372,0.718) (0.412,0.681) (0.449,0.641) (0.483,0.599) (0.513,0.556) (0.540,0.512) (0.564,0.467) (0.588,0.410) (0.607,0.352) (0.622,0.294) (0.631,0.236) (0.635,0.179) (0.634,0.123) (0.629,0.069) (0.619,0.016) (0.606,-0.035) (0.588,-0.084) (0.567,-0.130) (0.538,-0.181) (0.504,-0.228) (0.467,-0.271) (0.427,-0.308) (0.384,-0.341) (0.340,-0.368) (0.294,-0.390) (0.239,-0.410) (0.184,-0.422) (0.130,-0.428) (0.077,-0.428) (0.026,-0.421) (-0.023,-0.409) (-0.074,-0.390) (-0.122,-0.364) (-0.163,-0.333) (-0.199,-0.298) (-0.233,-0.255) (-0.259,-0.210) (-0.276,-0.162) (-0.287,-0.110) (-0.289,-0.058) (-0.282,-0.004) (-0.265,0.045) (-0.239,0.092) (-0.205,0.131) (-0.163,0.163) (-0.117,0.184) (-0.068,0.195) (-0.016,0.193) (0.033,0.178) (0.076,0.151) (0.109,0.113) (0.128,0.065) (0.131,0.014) (0.114,-0.035) (0.078,-0.072) (0.030,-0.089) (-0.020,-0.079) (-0.054,-0.042) (-0.056,0.008) (-0.017,0.040) (0.026,0.014)};
        \draw[wosm-accent, -latex] (-0.246,0.940) -- (-0.180,0.939);
        \end{scope}
      \end{tikzpicture}
      \par\small \textbf{Stable spiral}: \\ $\lambda = -\frac14 \pm i$. \\
      $A = \begin{pmatrix} -1/4 & 1 \\ -1 & -1/4 \end{pmatrix}$
    \end{column}
  \end{columns}
  \vskip0.4em
  {\small Dashed: the eigenvector lines, along which solutions move straight.\par}
\end{frame}

\begin{frame}{Classification by trace and determinant}
  For a $2 \times 2$ matrix, $\lambda_1 + \lambda_2 = \operatorname{tr} A = T$ and $\lambda_1 \lambda_2 = \det A = D$, so
  $\lambda = \frac12\bigl(T \pm \sqrt{T^2 - 4D}\bigr)$.
  \pause
  \begin{center}
    \begin{tabular}{lll}
      \toprule
      Condition & Eigenvalues & Origin \\
      \midrule
      $D < 0$ & real, opposite signs & saddle \\
      $D > 0$, \ $T^2 \ge 4D$ & real, same sign & node \\
      $D > 0$, \ $T^2 < 4D$, \ $T \ne 0$ & complex, $\operatorname{Re}\lambda = T/2$ & spiral \\
      $D > 0$, \ $T = 0$ & purely imaginary & centre (closed orbits) \\
      \bottomrule
    \end{tabular}
  \end{center}
  \pause
  Nodes and spirals are \textbf{stable} (attracting) if $T < 0$ and unstable if $T > 0$.
  \pause
  {\small Example: the damped spring $\vx' = \begin{pmatrix} 0 & 1 \\ -k/m & -c/m \end{pmatrix}\vx$ has $T = -c/m < 0$ and
  $D = k/m > 0$: stable, a spiral exactly when underdamped.\par}
\end{frame}

\begin{frame}{Stability}
  \begin{theorem}
    Every solution of $\vx' = A\vx$ tends to $\mathbf 0$ as $t \to \infty$ if and only if every eigenvalue of
    $A$ has negative real part.
  \end{theorem}
  {\small (For diagonalisable $A$: from $\sum_i c_i e^{\lambda_i t} \vv_i$ and $|e^{\lambda t}| = e^{(\operatorname{Re}\lambda)\, t}$.)\par}
  \pause
  \begin{block}{Linearisation (outlook)}
    Near an equilibrium $\vx^*$ of a nonlinear system $\vx' = F(\vx)$, the Jacobian $DF(\vx^*)$ (Lecture 12)
    usually decides stability.
  \end{block}
  \pause
  {\small Example: the gradient flow $\vx' = -\nabla f(\vx)$ has Jacobian $-H_f$. At a strict local minimum
  $H_f$ is positive definite (Lecture 14), so the minimum attracts. Gradient descent is Euler's method
  for this flow.\par}
\end{frame}

% =============================================================
\section{The matrix exponential}

\begin{frame}{The matrix exponential}
  \begin{definition}
    For a square matrix $A$: \quad $\displaystyle e^{tA} = I + tA + \frac{t^2 A^2}{2!} + \frac{t^3 A^3}{3!} + \cdots$
  \end{definition}
  \pause
  \begin{itemize}
    \item The series converges for every $A$ and $t$: each entry of $t^k A^k / k!$ is at most
      $(|t|\, n M)^k / k!$ in absolute value, where $M$ is the largest $|a_{ij}|$ --- compare with the
      series of $e^{|t| n M}$ (Lectures 2 and 5).
    \pause
    \item Differentiating term by term (Lecture 5): $\frac{d}{dt} e^{tA} = A e^{tA}$. So
      \[ \vx(t) = e^{tA}\, \vx_0 \]
      solves $\vx' = A\vx$, $\vx(0) = \vx_0$ --- exactly like $x(t) = e^{at} x_0$ in one variable.
  \end{itemize}
\end{frame}

\begin{frame}{Computing $e^{tA}$}
  \begin{itemize}
    \item \textbf{Diagonalisable}: $A = V\Lambda V^{-1}$ gives $A^k = V\Lambda^k V^{-1}$, so
      \[ e^{tA} = V\, \operatorname{diag}\bigl(e^{\lambda_1 t}, \dots, e^{\lambda_n t}\bigr)\, V^{-1} . \]
      Then $e^{tA}\vx_0$ is the eigenvector solution from before.
    \pause
    \item \textbf{A rotation}: $A = \begin{pmatrix} 0 & 1 \\ -1 & 0 \end{pmatrix}$ has $A^2 = -I$. Separating even and
      odd powers, as for Euler's formula:
      \[ e^{tA} = I\cos t + A \sin t = \begin{pmatrix} \cos t & \sin t \\ -\sin t & \cos t \end{pmatrix} . \]
  \end{itemize}
\end{frame}

\begin{frame}{When $A$ is not diagonalisable}
  $A = \begin{pmatrix} \lambda & 1 \\ 0 & \lambda \end{pmatrix} = \lambda I + N$, with
  $N = \begin{pmatrix} 0 & 1 \\ 0 & 0 \end{pmatrix}$ and $N^2 = 0$, has only one eigenvector direction.
  \pause
  \vskip0.4em
  Since $\lambda I$ and $N$ commute, $e^{t(\lambda I + N)} = e^{\lambda t}\, e^{tN}$, and the series for $e^{tN}$ stops
  after two terms:
  \[ e^{tA} = e^{\lambda t}\, (I + tN) = e^{\lambda t}\begin{pmatrix} 1 & t \\ 0 & 1 \end{pmatrix} . \]
  \pause
  {\small The factor $t e^{\lambda t}$ is the double root of Lecture 24 again. In general every matrix has a
  Jordan form (Linear Algebra, Unit 24), and $e^{tA}$ contains terms $t^k e^{\lambda t}$.\par}
\end{frame}

% =============================================================
\section{Looking back}

\begin{frame}{The module in one slide}
  \begin{itemize}
    \item \textbf{Series} (Units 1--8): convergence tests, power and Taylor series, generating functions,
      Stirling's formula.
    \item \textbf{Differential calculus in several variables} (Units 9--17): gradients, the chain rule
      and Jacobians, the Hessian and Taylor's formula, optimisation with and without constraints,
      gradient descent, implicit and inverse functions.
    \item \textbf{Integration} (Units 18--22): double and triple integrals, change of variables, the
      Gaussian integral, line integrals.
    \item \textbf{Differential equations} (Units 23--25): first-order, second-order linear, and linear
      systems.
  \end{itemize}
  \pause
  {\small Some threads: power series gave $e^{i\theta}$ and $e^{tA}$; eigenvalues decided the Hessian test,
  the speed of gradient descent, Lagrange problems on the sphere, and stability today; the Jacobian
  determinant governed the chain rule, inverse functions and the change of variables.\par}
\end{frame}

% =============================================================
\section{Summary and practice}

\begin{frame}{Summary}
  \begin{itemize}
    \item $\vx' = A\vx$: unique solutions for all $t$; an $n$-dimensional space of solutions.
    \item Eigenpairs give solutions $e^{\lambda t}\vv$; with $n$ independent eigenvectors,
      $\vx = \sum c_i e^{\lambda_i t}\vv_i$. Complex eigenvalues: take real and imaginary parts.
    \item $2 \times 2$ phase portraits: saddle, node, spiral, centre; classified by $\operatorname{tr} A$ and
      $\det A$. Stable iff all eigenvalues have negative real part.
    \item The matrix exponential: $\vx(t) = e^{tA}\vx_0$, with $e^{tA} = V e^{t\Lambda} V^{-1}$ when $A$ is
      diagonalisable.
  \end{itemize}
\end{frame}

\begin{frame}{Practice}
  \small
  \begin{enumerate}
    \item Solve $\vx' = A\vx$ with $A = \begin{pmatrix} 4 & -2 \\ 1 & 1 \end{pmatrix}$ and $\vx(0) = (3, 2)$.
    \item Classify the origin for $\begin{pmatrix} 1 & 2 \\ 3 & 2 \end{pmatrix}$, \
      $\begin{pmatrix} -3 & 1 \\ 1 & -3 \end{pmatrix}$, \ $\begin{pmatrix} 1 & -5 \\ 1 & -1 \end{pmatrix}$.
    \item Solve $x' = y$, $y' = -4x$ with $x(0) = 1$, $y(0) = 0$. What curve does $(x(t), y(t))$ trace?
    \item Write $y'' + 3y' + 2y = 0$ as a system, and compare the eigenvalues of its matrix with the
      roots of the characteristic equation.
    \item Compute $e^{tA}$ for $A = \begin{pmatrix} 0 & 1 \\ 0 & 0 \end{pmatrix}$, and solve $\vx' = A\vx$.
    \item For which $k$ do all solutions of $\vx' = \begin{pmatrix} -1 & k \\ 1 & -1 \end{pmatrix}\vx$ tend to $\mathbf 0$?
    \item Show directly that if $A\vv = \lambda\vv$, then $e^{\lambda t}\vv$ solves $\vx' = A\vx$.
  \end{enumerate}
\end{frame}

\begin{frame}{Answers}
  \small
  \begin{enumerate}
    \item $\lambda = 2$, $\vv = (1, 1)$; \ $\lambda = 3$, $\vv = (2, 1)$; \ $(3, 2) = (1, 1) + (2, 1)$: \
      $\vx = e^{2t}(1, 1) + e^{3t}(2, 1)$.
    \item $D = -4$: saddle. \ $T = -6$, $D = 8$, $T^2 > 4D$: stable node. \ $T = 0$, $D = 4$: centre.
    \item $x = \cos 2t$, $y = -2\sin 2t$: the ellipse $x^2 + \frac{y^2}{4} = 1$, traversed clockwise.
    \item $\begin{pmatrix} 0 & 1 \\ -2 & -3 \end{pmatrix}$: $\lambda^2 + 3\lambda + 2 = 0$, so $\lambda = -1, -2$ ---
      the same as the roots of $r^2 + 3r + 2 = 0$.
    \item $A^2 = 0$: \ $e^{tA} = I + tA = \begin{pmatrix} 1 & t \\ 0 & 1 \end{pmatrix}$; \
      $x_1 = x_1(0) + t\, x_2(0)$, $x_2 = x_2(0)$.
    \item $T = -2 < 0$; need $D = 1 - k > 0$: \ $k < 1$.
    \item $(e^{\lambda t}\vv)' = \lambda e^{\lambda t}\vv = e^{\lambda t} A\vv = A(e^{\lambda t}\vv)$.
  \end{enumerate}
\end{frame}

\begin{frame}{The end of the module}
  This was the last unit of Calculus~II.
  \begin{itemize}
    \item \textbf{Final exam}: written, $3$ hours, at the end of term ($45\%$ of the module). It tests
      computation (series, derivatives and integrals in several variables), the choice of method, and
      short proofs about convergence.
    \item \textbf{What uses this module}: multiple integrals and the Gaussian integral in Probability
      Theory and Statistics (joint densities); series and generating functions in the analysis of
      algorithms; gradients, Hessians and gradient descent in optimisation and machine learning.
  \end{itemize}
  \vskip1em
  {\small
  Questions: the module channel on Discord, or office hours,
  Monday--Friday, 10:00--17:00 (Baku time, UTC+4),
  \href{mailto:ali.khudiyev@wosm.academy}{ali.khudiyev@wosm.academy}.\par}
\end{frame}

\end{document}
